\section{Introduction}
In many applications it becomes necessary to evaluate
\[
   P = a_0 + a_1 x + a_2 x^2 + \dots + a_{n-1} x^{n-1} + x^{n},
\]
where the coefficients $a_i$ are known to the programmer~\cite{knuth1962evaluation}.
Horner's method evaluates $P$ using just $n$ multiplications:
$P = a_0 + x(a_1 + x(a_2 + \dots + x(a_{n-1} + x)\dots))$ and has been known for centuries.\footnote{Although named after William George Horner, this method is much older, as it has been attributed to Joseph-Louis Lagrange by Horner himself, and can be traced back many hundreds of years to Chinese and Persian mathematicians; see, e.g., Knuth~\cite{knuth1962evaluation}.}
Perhaps surprisingly, this is not optimal.

The trick is to allow some preprocessing of the coefficients $a_0, \dots, a_n$.
For example, the polynomial $9 + 7 x + 5 x^2 + 3 x^3 + x^4$ can be evaluated as
$y = x (x + 1)$, $P = (x + y - 1)(y + 4) + 13$ using just two multiplications rather than four.
In the 50s Motzkin gave an algorithm to evaluate any monic polynomial (i.e. leading coefficient 1) of degree $n$ using just $n/2$ multiplications.
%Later Kedlaya and Umanski~\cite{kedlaya2011fast} reduced this to just $(\log n)^{O(1)}$ time, using table lookups and advanced algebra.
The downside is that the method requires solving a polynomial system of equations over the complex numbers, which reduces precision and breaks over finite fields.
%Kedlaya and Umanski's method has a hidden factor of $\log_2 q$, where $q$ is the field size, making it impractical for low degree polynomials over large fields.
We will discuss the large body of related work by Paterson and Stockmeyer~\cite{paterson1973evaluation}, Kedlaya and Umanski~\cite{kedlaya2011fast}, and others in \cref{sec:related}.
Arguably, the most practical method is due to Rabin and Winograd~\cite{rabin1972number}, who gave a way to preprocess the coefficients using only rational operations, at the cost of $n/2 + 2\lceil\log n\rceil$ multiplications.
They conjectured this was optimal up to the constant 2.

In this paper we answer Rabin and Winograd's conjecture in the negative by showing how to evaluate any monic polynomial of degree $n$ using just $n/2+1$ multiplications with \emph{rational preprocessing}, and we show this is optimal.
At a high level, rational preprocessing is exactly what turns ``fast evaluation'' into a \emph{bijective/injective} parameterization: our evaluation circuit is parameterized by some $\theta$, and it induces a map from parameters $\theta$ to the coefficient vector $(a_0,\dots,a_{n-1})$ of the resulting polynomial.
Admitting a rational preprocessing/decoding map means that this parameter-to-coefficient map has a rational inverse on a dense set, i.e.\ it is generically one-to-one (birational onto its image).

Fast evaluation of small low-degree polynomials is important in many applications, such as cryptography~\cite{DBLP:conf/fse/CarletGPQR12} and $k$-wise independent hashing~\cite{wegman1981new, patracscu2012power}.
We show that our method can be used to speed up both of these applications substantially.
In Section~\ref{sec:experiments}, we demonstrate speedups of $1.3\times$ to $1.9\times$ over standard Horner evaluation for $k$-wise independent hashing with $k \in \{5, \ldots, 9\}$ on modern ARM processors.

\paragraph{Main Results.}
Section~\ref{sec:model} formalizes our model (arithmetic circuits / straight-line programs) and the notion of rational preprocessing/decoding.
Our main results are:
\begin{itemize}
   \item \textbf{Upper bound.} For every $n\ge 1$ and every field of characteristic $\neq 2$, there is an evaluation scheme for monic degree-$n$ polynomials that uses $\lfloor n/2\rfloor+1$ multiplications and admits rational preprocessing (Theorem~\ref{thm:main}).
   For characteristic~2, we have explicit constructions up to degree~15 and conjecture the bound holds in general.
   \item \textbf{Lower bound.} We prove that degree-$6$ evaluation with rational decoding requires at least $4$ multiplications (Section~\ref{sec:lower}), and conjecture that degree-$2n$ requires at least $n+1$ multiplications in general.
   \item \textbf{Injective Polynomials}. We give an ``injective'' polynomial hashing construction requiring only a single random key (Section~\ref{sec:injective}) and application-level benchmarks for hashing, data structures, and prime-field workloads (Section~\ref{sec:experiments}).
\end{itemize}
We also explore concrete constructions in characteristic~2 (currently up to degree~15) in the accompanying code.
These are of special interest since they can be used in cryptographic applications such as block ciphers and pseudorandom generators.
